% !TEX program = xelatex
% Options for packages loaded elsewhere
\PassOptionsToPackage{unicode}{hyperref}
\PassOptionsToPackage{hyphens}{url}
\PassOptionsToPackage{dvipsnames,svgnames,x11names}{xcolor}
\documentclass[
  8pt,
  twocolumn]{extarticle}
\usepackage{xcolor}
\usepackage[landscape,margin=0.38in]{geometry}
\usepackage{amsmath,amssymb}
\setcounter{secnumdepth}{-\maxdimen} % remove section numbering
\usepackage{iftex}
\ifPDFTeX
  \usepackage[T1]{fontenc}
  \usepackage[utf8]{inputenc}
  \usepackage{textcomp} % provide euro and other symbols
\else % if luatex or xetex
  \usepackage{unicode-math} % this also loads fontspec
  \defaultfontfeatures{Scale=MatchLowercase}
  \defaultfontfeatures[\rmfamily]{Ligatures=TeX,Scale=1}
\fi
\usepackage{lmodern}
\ifPDFTeX\else
  % xetex/luatex font selection
\fi
% Use upquote if available, for straight quotes in verbatim environments
\IfFileExists{upquote.sty}{\usepackage{upquote}}{}
\IfFileExists{microtype.sty}{% use microtype if available
  \usepackage[]{microtype}
  \UseMicrotypeSet[protrusion]{basicmath} % disable protrusion for tt fonts
}{}
\makeatletter
\@ifundefined{KOMAClassName}{% if non-KOMA class
  \IfFileExists{parskip.sty}{%
    \usepackage{parskip}
  }{% else
    \setlength{\parindent}{0pt}
    \setlength{\parskip}{6pt plus 2pt minus 1pt}}
}{% if KOMA class
  \KOMAoptions{parskip=half}}
\makeatother
\setlength{\emergencystretch}{3em} % prevent overfull lines
\providecommand{\tightlist}{%
  \setlength{\itemsep}{0pt}\setlength{\parskip}{0pt}}
\usepackage{microtype}
\usepackage{mathtools}
\usepackage{amssymb}
\setlength{\columnsep}{0.32in}
\setlength{\parindent}{0pt}
\setlength{\parskip}{2pt}
\setlength{\abovedisplayskip}{4pt}
\setlength{\belowdisplayskip}{4pt}
\setlength{\abovedisplayshortskip}{3pt}
\setlength{\belowdisplayshortskip}{3pt}
\allowdisplaybreaks
\sloppy
\emergencystretch=3em
\usepackage{bookmark}
\IfFileExists{xurl.sty}{\usepackage{xurl}}{} % add URL line breaks if available
\urlstyle{same}
\hypersetup{
  colorlinks=true,
  linkcolor={black},
  filecolor={Maroon},
  citecolor={Blue},
  urlcolor={black},
  pdfcreator={LaTeX via pandoc}}

\author{}
\date{}

\begin{document}

\section{All Formulas From Chatnotes}\label{all-formulas-from-chatnotes}

This sheet consolidates the formulas, model equations, hypothesis tests,
and exam templates contained in the \texttt{chatnotes} folder. Repeated
worked-example substitutions are consolidated into their underlying
formula so the sheet is usable for PDF export.

\subsection{Conventions}\label{conventions}

\[
t,T,h,k,n,p,q,N
\]

Variables: \(t\) is a time index; \(T\) is sample size or the final
observed time; \(h\) is a forecast horizon; \(k\) is a lag or number of
model parameters depending on context; \(n\) is a return horizon;
\(p,q\) are lag orders; \(N\) is the number of assets.

\[
F_t
\]

Variables: \(F_t\) is the information set available at time \(t\).

Returns may be written as decimals or percentages. Check the scale
before applying VaR or standard deviation formulas.

\subsection{Prices, Returns, and VaR}\label{prices-returns-and-var}

\subsubsection{Simple return}\label{simple-return}

\[
R_t=\frac{P_t-P_{t-1}}{P_{t-1}}=\frac{P_t}{P_{t-1}}-1
\]

Variables: \(R_t\) is the simple return from \(t-1\) to \(t\); \(P_t\)
is the price at time \(t\); \(P_{t-1}\) is the previous price.

\subsubsection{Gross return}\label{gross-return}

\[
1+R_t=\frac{P_t}{P_{t-1}}
\]

Variables: \(1+R_t\) is the gross return, the multiplier applied to
wealth.

\subsubsection{Return with dividends}\label{return-with-dividends}

\[
R_t=\frac{P_t+D_t-P_{t-1}}{P_{t-1}}
\]

Variables: \(D_t\) is the dividend paid during period \(t\).

\subsubsection{Log return}\label{log-return}

\[
r_t=\log(1+R_t)=\log(P_t)-\log(P_{t-1})=\log\left(\frac{P_t}{P_{t-1}}\right)
\]

Variables: \(r_t\) is the log return; \(\log\) is the natural logarithm.

\subsubsection{Small-return
approximation}\label{small-return-approximation}

\[
\log(1+R_t)\approx R_t
\]

Variables: \(R_t\) is the simple return; the approximation works best
when \(R_t\) is small.

\subsubsection{Multi-period gross
return}\label{multi-period-gross-return}

\[
1+R_t(k)=\prod_{j=0}^{k-1}(1+R_{t-j})
\]

Variables: \(R_t(k)\) is the \(k\)-period simple return ending at \(t\);
\(j\) indexes the returns being multiplied.

\subsubsection{Multi-period simple
return}\label{multi-period-simple-return}

\[
R_t(k)=\prod_{j=0}^{k-1}(1+R_{t-j})-1
\]

Variables: \(R_t(k)\) is the total simple return over \(k\) periods.

\subsubsection{Equity curve for one
dollar}\label{equity-curve-for-one-dollar}

\[
V_t(k)=\prod_{j=0}^{k-1}(1+R_{t-j})
\]

Variables: \(V_t(k)\) is the value at time \(t\) of one dollar invested
over \(k\) periods.

\subsubsection{Wealth growth}\label{wealth-growth}

\[
V_T=V_0\prod_{t=1}^{T}(1+R_t)
\]

Variables: \(V_0\) is starting wealth; \(V_T\) is ending wealth after
\(T\) periods.

\subsubsection{Multi-period log return}\label{multi-period-log-return}

\[
r_t(k)=\sum_{j=0}^{k-1}r_{t-j}=r_t+r_{t-1}+\cdots+r_{t-k+1}
\]

Variables: \(r_t(k)\) is the \(k\)-period log return.

\subsubsection{Total log return between two
prices}\label{total-log-return-between-two-prices}

\[
r(1,T)=\sum_{t=1}^{T}r_t=\log\left(\frac{P_T}{P_0}\right)
\]

Variables: \(P_0\) is the starting price; \(P_T\) is the ending price.

\subsubsection{Index log return}\label{index-log-return}

\[
r_t=\log(Index_t)-\log(Index_{t-1})
\]

Variables: \(Index_t\) is the index level at time \(t\).

\subsubsection{Excess returns}\label{excess-returns}

\[
\begin{aligned}
\text{asset excess return} &= \text{asset return}-R_f \\
\text{market excess return} &= \text{market return}-R_f
\end{aligned}
\]

Variables: \(R_f\) is the risk-free return.

\subsubsection{Market beta
interpretation}\label{market-beta-interpretation}

\[
\begin{aligned}
\beta>1 &\Rightarrow \text{aggressive, more sensitive than the market} \\
\beta=1 &\Rightarrow \text{tracks the market} \\
0<\beta<1 &\Rightarrow \text{defensive, less sensitive than the market} \\
\beta=0 &\Rightarrow \text{unrelated to market movements}
\end{aligned}
\]

Variables: \(\beta\) is market sensitivity.

\subsubsection{Market capitalisation}\label{market-capitalisation}

\[
\text{market capitalisation}=\text{share price}\times\text{number of shares outstanding}
\]

Variables: market capitalisation is firm size measured using market
value.

\subsubsection{Book-to-market ratio}\label{book-to-market-ratio}

\[
\operatorname{BtM}=\frac{\text{book value of equity}}{\text{market value of equity}}
\]

Variables: \(\operatorname{BtM}\) is the book-to-market ratio.

\subsubsection{Normal return model}\label{normal-return-model}

\[
R\sim\mathcal N(\mu,\sigma^2)
\]

Variables: \(R\) is a return; \(\mu\) is mean return; \(\sigma^2\) is
variance; \(\sigma\) is standard deviation.

\subsubsection{Normal quantile}\label{normal-quantile}

\[
q_\alpha=\mu+\sigma z_\alpha
\]

Variables: \(q_\alpha\) is the lower-tail return quantile; \(\alpha\) is
the tail probability; \(z_\alpha\) is the standard normal cutoff.

\subsubsection{Common normal cutoffs}\label{common-normal-cutoffs}

\[
\begin{aligned}
z_{0.05}&=-1.645\\
z_{0.025}&=-1.96
\end{aligned}
\]

Variables: \(z_{0.05}\) is the 5 percent lower-tail cutoff;
\(z_{0.025}\) is the 2.5 percent lower-tail cutoff.

\subsubsection{One-asset VaR}\label{one-asset-var}

\[
\operatorname{VaR}_\alpha=|W_0q_\alpha|
\]

Variables: \(\operatorname{VaR}_\alpha\) is Value at Risk at tail
probability \(\alpha\); \(W_0\) is the dollar value invested.

\subsubsection{Zero-mean 5 percent normal
VaR}\label{zero-mean-5-percent-normal-var}

\[
\begin{aligned}
q_{0.05}&=-1.645\sigma\\
\operatorname{VaR}_{0.05}&=|W_0(-1.645\sigma)|
\end{aligned}
\]

Variables: \(q_{0.05}\) is the 5 percent return cutoff; \(\sigma\) is
return volatility.

\subsection{Random Variables and Distribution
Properties}\label{random-variables-and-distribution-properties}

\subsubsection{Probability under a
density}\label{probability-under-a-density}

\[
\Pr(a\le X\le b)=\text{area under }f(x)\text{ from }a\text{ to }b
\]

Variables: \(X\) is a random variable; \(a,b\) are cutoffs; \(f(x)\) is
the probability density.

\subsubsection{Normal distribution
notation}\label{normal-distribution-notation}

\[
X\sim\mathcal N(\mu,\sigma^2)
\]

Variables: \(X\) is a random variable; \(\mu\) is the mean; \(\sigma^2\)
is the variance.

\subsubsection{Expected value}\label{expected-value}

\[
\mathbb E[X]=\mu
\]

Variables: \(\mathbb E[\cdot]\) means expected value.

\subsubsection{State-contingent expected
return}\label{state-contingent-expected-return}

\[
\mathbb E(R)=\sum_s R_s\Pr(S=s)
\]

Variables: \(R_s\) is the return in state \(s\); \(S\) is the state
variable.

\subsubsection{Expected value linearity}\label{expected-value-linearity}

\[
\begin{aligned}
\mathbb E(a)&=a\\
\mathbb E(aX)&=a\mathbb E(X)\\
\mathbb E(a+X)&=a+\mathbb E(X)\\
\mathbb E(X+Y)&=\mathbb E(X)+\mathbb E(Y)
\end{aligned}
\]

Variables: \(a\) is a constant; \(X,Y\) are random variables.

\subsubsection{Variance}\label{variance}

\[
\operatorname{Var}(X)=\mathbb E[(X-\mu)^2]
\]

Variables: \(\operatorname{Var}(X)\) is the variance of \(X\).

\subsubsection{Standard deviation}\label{standard-deviation}

\[
\operatorname{sd}(X)=\sqrt{\operatorname{Var}(X)}
\]

Variables: \(\operatorname{sd}(X)\) is the standard deviation of \(X\).

\subsubsection{Variance rules}\label{variance-rules}

\[
\begin{aligned}
\operatorname{Var}(a)&=0\\
\operatorname{Var}(aX)&=a^2\operatorname{Var}(X)\\
\operatorname{Var}(a+X)&=\operatorname{Var}(X)
\end{aligned}
\]

Variables: \(a\) is a constant.

\subsubsection{Variance of a sum}\label{variance-of-a-sum}

\[
\operatorname{Var}(X+Y)=\operatorname{Var}(X)+\operatorname{Var}(Y)+2\operatorname{Cov}(X,Y)
\]

Variables: \(\operatorname{Cov}(X,Y)\) is covariance between \(X\) and
\(Y\).

\subsubsection{Standardising a normal
variable}\label{standardising-a-normal-variable}

\[
Z=\frac{X-\mu}{\sigma}\sim\mathcal N(0,1)
\]

Variables: \(Z\) is a standard normal variable; \(\sigma\) is the
standard deviation of \(X\).

\subsubsection{Probability conversion after
standardising}\label{probability-conversion-after-standardising}

\[
\Pr(X>a)=\Pr\left(\frac{X-\mu}{\sigma}>\frac{a-\mu}{\sigma}\right)=\Pr(Z>z)
\]

Variables: \(z=(a-\mu)/\sigma\) is the standardised cutoff.

\subsubsection{Lower-tail normal
probability}\label{lower-tail-normal-probability}

\[
\Pr(r_t<x)=\Pr\left(Z<\frac{x-\mu}{\sigma}\right)
\]

Variables: \(r_t\) is a return; \(x\) is a cutoff.

\subsubsection{Return scale convention}\label{return-scale-convention}

\[
\begin{aligned}
0.01&=1\text{ percent, if returns are decimals}\\
1&=1\text{ percent, if returns are stored as percentage points}
\end{aligned}
\]

Variables: the same economic return can be represented on different
numeric scales.

\subsubsection{Sample mean}\label{sample-mean}

\[
\bar r=\frac{1}{T}\sum_{t=1}^{T}r_t
\]

Variables: \(\bar r\) is the sample average return.

\subsubsection{Sample variance and standard
deviation}\label{sample-variance-and-standard-deviation}

\[
s^2=\frac{1}{T-1}\sum_{t=1}^{T}(r_t-\bar r)^2,\qquad s=\sqrt{s^2}
\]

Variables: \(s^2\) is sample variance; \(s\) is sample standard
deviation.

\subsubsection{Skewness and kurtosis}\label{skewness-and-kurtosis}

\[
\text{skewness}=\frac{\mathbb E[(r_t-\mu)^3]}{\sigma^3},\qquad
\text{kurtosis}=\frac{\mathbb E[(r_t-\mu)^4]}{\sigma^4}
\]

Variables: skewness measures asymmetry; kurtosis measures tail
thickness.

\subsection{Co-Movement, Conditional Moments, and
Factors}\label{co-movement-conditional-moments-and-factors}

\subsubsection{Joint probability}\label{joint-probability}

\[
\Pr(X=x,Y=y)
\]

Variables: \(X,Y\) are random variables; \(x,y\) are outcomes.

\subsubsection{Marginal probability}\label{marginal-probability}

\[
\Pr(X=x)=\sum_y\Pr(X=x,Y=y)
\]

Variables: the sum is over all possible values of \(Y\).

\subsubsection{Conditional probability}\label{conditional-probability}

\[
\Pr(X=x\mid Y=y)=\frac{\Pr(X=x,Y=y)}{\Pr(Y=y)}
\]

Variables: \(\Pr(X=x\mid Y=y)\) is the probability of \(X=x\) given
\(Y=y\).

\subsubsection{Independence}\label{independence}

\[
\Pr(X=x\mid Y=y)=\Pr(X=x)
\]

Variables: independence means knowing \(Y\) does not change the
probability distribution of \(X\).

\subsubsection{Conditional expectation}\label{conditional-expectation}

\[
\mathbb E[X\mid Y]
\]

Variables: \(\mathbb E[X\mid Y]\) is the expected value of \(X\)
conditional on \(Y\).

\subsubsection{Law of iterated
expectations}\label{law-of-iterated-expectations}

\[
\mathbb E[X]=\mathbb E[\mathbb E[X\mid Y]]
\]

Variables: \(Y\) is the conditioning variable.

\subsubsection{Conditional variance}\label{conditional-variance}

\[
\operatorname{Var}(Y\mid X)=\mathbb E[(Y-\mathbb E[Y\mid X])^2\mid X]
\]

Variables: \(\operatorname{Var}(Y\mid X)\) is remaining uncertainty
about \(Y\) after conditioning on \(X\).

\subsubsection{Law of total variance}\label{law-of-total-variance}

\[
\operatorname{Var}(Y)=\mathbb E[\operatorname{Var}(Y\mid X)]+\operatorname{Var}(\mathbb E[Y\mid X])
\]

Variables: the first term is average conditional variance; the second is
variance of conditional means.

\subsubsection{Information reduces conditional
variance}\label{information-reduces-conditional-variance}

\[
\operatorname{Var}(Y\mid \text{information})<\operatorname{Var}(Y)
\]

Variables: information is any conditioning set that improves prediction.

\subsubsection{Covariance}\label{covariance}

\[
\operatorname{Cov}(X,Y)=\mathbb E[(X-\mathbb E[X])(Y-\mathbb E[Y])]
\]

Variables: covariance measures linear co-movement.

\subsubsection{State covariance}\label{state-covariance}

\[
\operatorname{Cov}(R_A,R_B)=\sum_s p_s(R_{A,s}-\mu_A)(R_{B,s}-\mu_B)
\]

Variables: \(R_A,R_B\) are asset returns; \(p_s\) is the probability of
state \(s\); \(\mu_A,\mu_B\) are expected returns.

\subsubsection{Covariance from
correlation}\label{covariance-from-correlation}

\[
\operatorname{Cov}(R_A,R_B)=\rho_{AB}\sigma_A\sigma_B
\]

Variables: \(\rho_{AB}\) is the correlation between assets A and B;
\(\sigma_A,\sigma_B\) are standard deviations.

\subsubsection{Correlation}\label{correlation}

\[
\operatorname{Corr}(X,Y)=\frac{\operatorname{Cov}(X,Y)}{\operatorname{sd}(X)\operatorname{sd}(Y)}
\]

Variables: \(\operatorname{Corr}(X,Y)\) is the standardised covariance.

\subsubsection{Correlation from
covariance}\label{correlation-from-covariance}

\[
\rho_{AB}=\frac{\operatorname{Cov}(R_A,R_B)}{\sigma_A\sigma_B}
\]

Variables: \(\rho_{AB}\) is the correlation coefficient.

\subsubsection{Covariance sign
interpretation}\label{covariance-sign-interpretation}

\[
\begin{aligned}
\operatorname{Cov}(X,Y)>0&\Rightarrow X\text{ and }Y\text{ tend to move together}\\
\operatorname{Cov}(X,Y)<0&\Rightarrow X\text{ and }Y\text{ tend to move oppositely}\\
\operatorname{Cov}(X,Y)=0&\Rightarrow \text{no linear co-movement}
\end{aligned}
\]

Variables: the sign of covariance gives direction of linear co-movement.

\subsubsection{Group mean return}\label{group-mean-return}

\[
\bar r_g=\frac{1}{T}\sum_{t=1}^{T}r_{g,t}
\]

Variables: \(\bar r_g\) is the mean return for group \(g\); \(r_{g,t}\)
is group return at time \(t\).

\subsubsection{Generic factor long-short
return}\label{generic-factor-long-short-return}

\[
r_{\text{factor},t}=r_{\text{long},t}-r_{\text{short},t}
\]

Variables: \(r_{\text{long},t}\) is the return on the long leg;
\(r_{\text{short},t}\) is the return on the short leg.

\subsubsection{SMB factor}\label{smb-factor}

\[
SMB_t=R_{\text{small},t}-R_{\text{big},t}
\]

Variables: \(SMB_t\) is Small Minus Big at time \(t\).

\subsubsection{HML factor}\label{hml-factor}

\[
HML_t=R_{\text{high BtM},t}-R_{\text{low BtM},t}
\]

Variables: \(HML_t\) is High Minus Low book-to-market at time \(t\).

\subsubsection{Fama-French two-by-three SMB
construction}\label{fama-french-two-by-three-smb-construction}

\[
SMB=\operatorname{average}(S/L,S/N,S/H)-\operatorname{average}(B/L,B/N,B/H)
\]

Variables: \(S\) means small; \(B\) means big; \(L,N,H\) mean low,
neutral, and high book-to-market portfolios.

\subsubsection{Fama-French two-by-three HML
construction}\label{fama-french-two-by-three-hml-construction}

\[
HML=\operatorname{average}(S/H,B/H)-\operatorname{average}(S/L,B/L)
\]

Variables: \(S/H\) is the small high-BtM portfolio; \(B/H\) is the big
high-BtM portfolio; \(S/L,B/L\) are low-BtM portfolios.

\subsection{Portfolio Risk and Sharpe
Ratios}\label{portfolio-risk-and-sharpe-ratios}

\subsubsection{Two-asset portfolio
return}\label{two-asset-portfolio-return}

\[
R_p=wR_1+(1-w)R_2
\]

Variables: \(R_p\) is portfolio return; \(w\) is the weight in asset 1;
\(R_1,R_2\) are asset returns.

\subsubsection{Two-asset expected
return}\label{two-asset-expected-return}

\[
\mathbb E[R_p]=w\mathbb E[R_1]+(1-w)\mathbb E[R_2]
\]

Variables: \(\mathbb E[R_p]\) is expected portfolio return.

\subsubsection{Two-asset expected return using
means}\label{two-asset-expected-return-using-means}

\[
\mu_p=w\mu_1+(1-w)\mu_2
\]

Variables: \(\mu_p\) is portfolio mean return; \(\mu_1,\mu_2\) are asset
mean returns.

\subsubsection{Two-asset portfolio
variance}\label{two-asset-portfolio-variance}

\[
\operatorname{Var}(R_p)=w^2\sigma_1^2+(1-w)^2\sigma_2^2+2w(1-w)\sigma_{12}
\]

Variables: \(\sigma_1^2,\sigma_2^2\) are asset variances;
\(\sigma_{12}\) is covariance.

\subsubsection{Two-asset portfolio variance using explicit
weights}\label{two-asset-portfolio-variance-using-explicit-weights}

\[
\sigma_p^2=w_1^2\sigma_1^2+w_2^2\sigma_2^2+2w_1w_2\sigma_{12}
\]

Variables: \(w_1,w_2\) are portfolio weights; \(\sigma_p^2\) is
portfolio variance.

\subsubsection{Covariance from
correlation}\label{covariance-from-correlation-1}

\[
\sigma_{12}=\rho_{12}\sigma_1\sigma_2
\]

Variables: \(\rho_{12}\) is correlation between assets 1 and 2.

\subsubsection{Sharpe ratio}\label{sharpe-ratio}

\[
SR=\frac{\operatorname{mean\ return}-\text{risk-free rate}}{\text{standard deviation}}
\]

Variables: \(SR\) is the Sharpe ratio.

\subsubsection{Asset Sharpe ratio}\label{asset-sharpe-ratio}

\[
SR_i=\frac{\mu_i-r_f}{\sigma_i}
\]

Variables: \(\mu_i\) is expected return on asset \(i\); \(r_f\) is the
risk-free rate; \(\sigma_i\) is standard deviation.

\subsubsection{Portfolio Sharpe ratio}\label{portfolio-sharpe-ratio}

\[
SR_p=\frac{\mu_p-r_f}{\sigma_p}
\]

Variables: \(SR_p\) is portfolio Sharpe ratio; \(\sigma_p\) is portfolio
standard deviation.

\subsubsection{Two-asset GMV weight}\label{two-asset-gmv-weight}

\[
w_1^{GMV}=\frac{\sigma_2^2-\sigma_{12}}{\sigma_1^2+\sigma_2^2-2\sigma_{12}}
\]

Variables: \(w_1^{GMV}\) is the global minimum variance weight in asset
1.

\subsubsection{Second GMV weight}\label{second-gmv-weight}

\[
w_2^{GMV}=1-w_1^{GMV}
\]

Variables: \(w_2^{GMV}\) is the GMV weight in asset 2.

\subsubsection{Regression trick for GMV
weights}\label{regression-trick-for-gmv-weights}

\[
r_{2,t}=\beta_0+\beta_1(r_{2,t}-r_{1,t})+e_t
\]

Variables: \(r_{1,t},r_{2,t}\) are asset returns; \(\beta_0,\beta_1\)
are regression coefficients; \(e_t\) is the error.

\[
\beta_1=w_1,\qquad w_2=1-\beta_1
\]

Variables: \(\beta_1\) gives the weight in asset 1 under the regression
trick.

\subsubsection{Regression slope form used in
practice}\label{regression-slope-form-used-in-practice}

\[
b=\frac{\operatorname{Cov}(R_1,R_2)}{\operatorname{Var}(R_2)}
\]

Variables: \(b\) is a regression slope in the practice note.

\subsubsection{Many-asset GMV weights}\label{many-asset-gmv-weights}

\[
\mathbf w^{GMV}=\frac{\Sigma^{-1}\mathbf 1}{\mathbf 1'\Sigma^{-1}\mathbf 1}
\]

Variables: \(\mathbf w^{GMV}\) is the GMV weight vector; \(\Sigma\) is
the covariance matrix; \(\mathbf 1\) is a vector of ones.

\subsubsection{Portfolio return from GMV
weights}\label{portfolio-return-from-gmv-weights}

\[
r_{p,t}=w_1r_{1,t}+(1-w_1)r_{2,t}
\]

Variables: \(r_{p,t}\) is portfolio return at time \(t\).

\subsection{Asset Pricing Models}\label{asset-pricing-models}

\subsubsection{CAPM pricing equation}\label{capm-pricing-equation}

\[
\mathbb E(R_i)-r_f=\beta_i[\mathbb E(R_m)-r_f]
\]

Variables: \(R_i\) is asset return; \(R_m\) is market return;
\(\beta_i\) is market beta.

\subsubsection{CAPM regression}\label{capm-regression}

\[
r_{i,t}-r_{f,t}=\alpha_i+\beta_i(r_{m,t}-r_{f,t})+e_{i,t}
\]

Variables: \(r_{i,t}-r_{f,t}\) is asset excess return;
\(r_{m,t}-r_{f,t}\) is market excess return; \(\alpha_i\) is abnormal
return; \(e_{i,t}\) is the residual.

\subsubsection{CAPM alpha restriction}\label{capm-alpha-restriction}

\[
\alpha_i=0
\]

Variables: \(\alpha_i\) is abnormal return relative to CAPM.

\subsubsection{CAPM total risk
decomposition}\label{capm-total-risk-decomposition}

\[
\operatorname{Var}(r_{i,t}-r_f)=\beta_i^2\operatorname{Var}(r_{m,t}-r_f)+\operatorname{Var}(e_{i,t})
\]

Variables: total excess-return variance equals systematic variance plus
idiosyncratic variance.

\subsubsection{Systematic and idiosyncratic
fractions}\label{systematic-and-idiosyncratic-fractions}

\[
\begin{aligned}
\text{systematic fraction}&=R^2\\
\text{idiosyncratic fraction}&=1-R^2
\end{aligned}
\]

Variables: \(R^2\) is the regression explained-variation statistic.

\subsubsection{Fama-French three-factor
model}\label{fama-french-three-factor-model}

\[
r_{i,t}-r_{f,t}=\alpha_i+\beta_{i1}(r_{m,t}-r_{f,t})+\beta_{i2}SMB_t+\beta_{i3}HML_t+e_{i,t}
\]

Variables: \(\beta_{i1}\) is market loading; \(\beta_{i2}\) is SMB
loading; \(\beta_{i3}\) is HML loading.

\subsubsection{Carhart four-factor
model}\label{carhart-four-factor-model}

\[
r_{i,t}-r_{f,t}=\alpha_i+\beta_{i1}(r_{m,t}-r_{f,t})+\beta_{i2}SMB_t+\beta_{i3}HML_t+\beta_{i4}MOM_t+e_{i,t}
\]

Variables: \(MOM_t\) is the momentum factor; \(\beta_{i4}\) is momentum
loading.

\subsubsection{CAPM nested in the four-factor
model}\label{capm-nested-in-the-four-factor-model}

\[
\beta_{i2}=\beta_{i3}=\beta_{i4}=0
\]

Variables: setting the extra factor loadings to zero reduces the
four-factor model to CAPM.

\subsubsection{Single-coefficient
t-test}\label{single-coefficient-t-test}

\[
\begin{aligned}
H_0&:\beta=\beta_0\\
H_1&:\beta\ne\beta_0\\
t&=\frac{\hat\beta-\beta_0}{\operatorname{se}(\hat\beta)}
\end{aligned}
\]

Variables: \(\hat\beta\) is the estimated coefficient; \(\beta_0\) is
the hypothesised value; \(\operatorname{se}(\hat\beta)\) is the standard
error.

\subsubsection{Test whether a factor
matters}\label{test-whether-a-factor-matters}

\[
\begin{aligned}
H_0&:\beta_{SMB}=0\\
H_1&:\beta_{SMB}\ne0\\
t&=\frac{\hat\beta_{SMB}}{\operatorname{se}(\hat\beta_{SMB})}
\end{aligned}
\]

Variables: \(\beta_{SMB}\) is the SMB coefficient.

\subsubsection{Joint test of CAPM against a multi-factor
model}\label{joint-test-of-capm-against-a-multi-factor-model}

\[
\begin{aligned}
H_0&:\beta_{SMB}=\beta_{HML}=\beta_{MOM}=0\\
H_1&:\text{at least one extra }\beta\ne0\\
J&=\frac{RSS_0-RSS_1}{RSS_1/(T-K-1)}
\end{aligned}
\]

Variables: \(RSS_0\) is restricted-model residual sum of squares;
\(RSS_1\) is unrestricted-model residual sum of squares; \(K\) is the
number of regressors in the unrestricted model.

\subsubsection{Residual standard error and
RSS}\label{residual-standard-error-and-rss}

\[
RSS=(\text{residual standard error})^2df
\]

Variables: \(df\) is residual degrees of freedom.

\[
\text{residual standard error}=\sqrt{\frac{RSS}{df}}
\]

Variables: residual standard error estimates residual standard
deviation.

\[
df=T-K-1
\]

Variables: \(T\) is sample size; \(K\) is the number of regressors
excluding the intercept.

\subsection{Regression Diagnostics and EMH
Tests}\label{regression-diagnostics-and-emh-tests}

\subsubsection{Regression residual
properties}\label{regression-residual-properties}

\[
\begin{aligned}
\mathbb E(e_{i,t})&=0\\
\operatorname{Var}(e_{i,t})&\text{ is constant}\\
\mathbb E(e_{i,t}e_{i,t-k})&=0,\quad k=1,2,\ldots
\end{aligned}
\]

Variables: \(e_{i,t}\) is a regression disturbance or residual.

\subsubsection{Original model and
residual}\label{original-model-and-residual}

\[
y_t=\alpha+\beta x_t+e_t
\]

Variables: \(y_t\) is the dependent variable; \(x_t\) is the regressor;
\(\alpha,\beta\) are coefficients.

\[
\hat e_t=y_t-\hat\alpha-\hat\beta x_t
\]

Variables: \(\hat e_t\) is the fitted residual.

\subsubsection{Heteroskedasticity
hypotheses}\label{heteroskedasticity-hypotheses}

\[
\begin{aligned}
H_0&:\operatorname{Var}(e_{i,t})\text{ is constant}\\
H_1&:\operatorname{Var}(e_{i,t})\text{ is not constant}
\end{aligned}
\]

Variables: the null is homoskedasticity; the alternative is
heteroskedasticity.

\subsubsection{Heteroskedasticity auxiliary
regression}\label{heteroskedasticity-auxiliary-regression}

\[
\hat e_t^2=\gamma_0+\gamma_1x_t+\gamma_2x_t^2+v_t
\]

Variables: \(\gamma_j\) are auxiliary-regression coefficients; \(v_t\)
is the auxiliary error.

\subsubsection{LM statistic for
heteroskedasticity}\label{lm-statistic-for-heteroskedasticity}

\[
W=TR^2
\]

Variables: \(R^2\) comes from the auxiliary regression; compare \(W\)
with a chi-square critical value.

\subsubsection{Residual autocorrelation
hypotheses}\label{residual-autocorrelation-hypotheses}

\[
\begin{aligned}
H_0&:\text{no autocorrelation}\\
H_1&:\text{at least one lagged residual matters}
\end{aligned}
\]

Variables: the alternative means residual serial correlation remains.

\subsubsection{Residual autocorrelation auxiliary
regression}\label{residual-autocorrelation-auxiliary-regression}

\[
\hat e_t=\gamma_0+\gamma_1\hat e_{t-1}+\gamma_2\hat e_{t-2}+\cdots+v_t
\]

Variables: \(\hat e_{t-j}\) are lagged residuals.

\subsubsection{LM statistic for
autocorrelation}\label{lm-statistic-for-autocorrelation}

\[
AR(p)=TR^2
\]

Variables: \(p\) is the number of residual lags; under the null compare
to \(\chi_p^2\).

\subsubsection{HAC definition}\label{hac-definition}

\[
HAC=\text{heteroskedasticity and autocorrelation consistent}
\]

Variables: HAC/Newey-West changes standard errors, not fitted
coefficients.

\subsubsection{ETF tracking regression}\label{etf-tracking-regression}

\[
r_{\text{ETF},t}=\alpha+\beta r_{\text{index},t}+e_t
\]

Variables: \(\beta\) measures the ETF's exposure to the index.

\subsubsection{Leveraged ETF target
test}\label{leveraged-etf-target-test}

\[
\begin{aligned}
H_0&:\beta=-2\\
H_1&:\beta\ne-2\\
t&=\frac{\hat\beta-(-2)}{\operatorname{se}(\hat\beta)}
\end{aligned}
\]

Variables: \(\hat\beta\) is the estimated tracking slope.

\subsubsection{EMH return model}\label{emh-return-model}

\[
r_t=\mu+e_t
\]

Variables: \(\mu\) is constant mean return; \(e_t\) is the unpredictable
return component.

\subsubsection{White noise notation}\label{white-noise-notation}

\[
e_t\sim WN(0,\sigma^2)
\]

Variables: \(WN(0,\sigma^2)\) means white noise with zero mean and
variance \(\sigma^2\).

\subsubsection{White noise conditions}\label{white-noise-conditions}

\[
\begin{aligned}
\mathbb E(e_t)&=0\\
\operatorname{Var}(e_t)&=\sigma^2\\
\operatorname{Corr}(e_t,e_{t-k})&=0,\quad k\ge1
\end{aligned}
\]

Variables: no autocorrelation means past shocks do not linearly predict
current shocks.

\subsubsection{Autocorrelation}\label{autocorrelation}

\[
\rho(k)=\operatorname{Corr}(r_t,r_{t-k})
\]

Variables: \(\rho(k)\) is autocorrelation at lag \(k\).

\subsubsection{Correlogram confidence
band}\label{correlogram-confidence-band}

\[
\pm\frac{2}{\sqrt T}
\]

Variables: \(T\) is sample size; the band is an approximate 5 percent
individual significance guide.

\subsubsection{Q-test hypotheses}\label{q-test-hypotheses}

\[
\begin{aligned}
H_0&:\rho_1=\rho_2=\cdots=\rho_k=0\\
H_1&:\text{at least one autocorrelation is non-zero}
\end{aligned}
\]

Variables: \(\rho_j\) is autocorrelation at lag \(j\).

\subsubsection{Portmanteau/Ljung-Box statistic from the practice
exam}\label{portmanteauljung-box-statistic-from-the-practice-exam}

\[
Q=T(T+2)\sum_{k=1}^{m}\frac{\hat\rho_k^2}{T-k}
\]

Variables: \(m\) is the largest lag tested; \(\hat\rho_k\) is the sample
autocorrelation at lag \(k\); compare \(Q\) with \(\chi_m^2\).

\subsubsection{Single autocorrelation large-sample
test}\label{single-autocorrelation-large-sample-test}

\[
\sqrt T\,\hat\rho(k)\approx\mathcal N(0,1)
\]

Variables: \(\hat\rho(k)\) is the sample autocorrelation at lag \(k\).

\[
\begin{aligned}
H_0&:\rho(k)=0\\
H_1&:\rho(k)\ne0\\
\left|\sqrt T\,\hat\rho(k)\right|&>1.96\Rightarrow \text{reject at approximately 5 percent}
\end{aligned}
\]

Variables: \(1.96\) is the two-sided 5 percent normal cutoff.

\subsubsection{Demeaned returns for LM predictability
tests}\label{demeaned-returns-for-lm-predictability-tests}

\[
e_t=r_t-\bar r
\]

Variables: \(\bar r\) is the sample mean return.

\subsubsection{LM return predictability
regression}\label{lm-return-predictability-regression}

\[
e_t=\gamma_0+\gamma_1e_{t-1}+\gamma_2e_{t-2}+\cdots+v_t
\]

Variables: lagged demeaned returns test predictability.

\subsubsection{LM return predictability hypotheses and
statistic}\label{lm-return-predictability-hypotheses-and-statistic}

\[
\begin{aligned}
H_0&:\gamma_1=\gamma_2=\cdots=\gamma_k=0\\
H_1&:\text{at least one }\gamma_j\text{ differs from zero}\\
AR(k)&=TR^2\sim\chi_k^2
\end{aligned}
\]

Variables: \(\gamma_j\) are lag coefficients; \(\chi_k^2\) is chi-square
with \(k\) degrees of freedom.

\subsubsection{Variance-ratio one-period log
return}\label{variance-ratio-one-period-log-return}

\[
r_t=\log(P_t)-\log(P_{t-1})
\]

Variables: \(r_t\) is one-period log return.

\subsubsection{Variance-ratio n-period log
return}\label{variance-ratio-n-period-log-return}

\[
r_t(n)=r_t+r_{t-1}+\cdots+r_{t-(n-1)}
\]

Variables: \(r_t(n)\) is the \(n\)-period log return.

\subsubsection{Variance scaling under
independence}\label{variance-scaling-under-independence}

\[
\operatorname{Var}[\text{n-period return}]=n\operatorname{Var}[\text{1-period return}]
\]

Variables: this holds approximately when returns are independent over
time.

\subsubsection{Variance ratio}\label{variance-ratio}

\[
VR_n=\frac{\text{variance of n-period returns}}{n\times\text{variance of 1-period returns}}
\]

Variables: \(VR_n=1\) suggests no autocorrelation; \(VR_n>1\) suggests
positive autocorrelation; \(VR_n<1\) suggests negative autocorrelation.

\subsection{Stationarity, Random Walks, ADF, and AR
Models}\label{stationarity-random-walks-adf-and-ar-models}

\subsubsection{Covariance stationarity}\label{covariance-stationarity}

\[
\begin{aligned}
\mathbb E(Y_t)&=\mu\\
\operatorname{Var}(Y_t)&=\sigma^2\\
\operatorname{Cov}(Y_t,Y_{t-j})&=\gamma_j
\end{aligned}
\]

Variables: stationarity means constant mean, constant variance, and
autocovariance depending only on lag \(j\).

\subsubsection{Stationarity rule of
thumb}\label{stationarity-rule-of-thumb}

\[
\begin{aligned}
\log\text{ prices}&\quad \text{often non-stationary}\\
\log\text{ returns}&\quad \text{often stationary}
\end{aligned}
\]

Variables: this is a modelling guide, not a formal test.

\subsubsection{Random walk}\label{random-walk}

\[
p_t=p_{t-1}+\epsilon_t
\]

Variables: \(p_t\) is usually log price; \(\epsilon_t\) is the shock.

\subsubsection{Random walk successive
substitution}\label{random-walk-successive-substitution}

\[
p_t=p_0+\epsilon_1+\cdots+\epsilon_t
\]

Variables: \(p_0\) is the initial value.

\subsubsection{Random walk moments}\label{random-walk-moments}

\[
\mathbb E(p_t)=p_0,\qquad \operatorname{Var}(p_t)=\sigma^2t
\]

Variables: \(\sigma^2\) is shock variance. If \(p_0=0\), the note's
simplified mean is \(\mathbb E(p_t)=0\).

\subsubsection{Random walk with drift}\label{random-walk-with-drift}

\[
p_t=\mu+p_{t-1}+\epsilon_t
\]

Variables: \(\mu\) is drift.

\subsubsection{Random walk with drift successive
substitution}\label{random-walk-with-drift-successive-substitution}

\[
p_t=p_0+\mu t+\epsilon_1+\cdots+\epsilon_t
\]

Variables: drift accumulates linearly over time.

\subsubsection{Random walk with drift
moments}\label{random-walk-with-drift-moments}

\[
\mathbb E(p_t)=p_0+\mu t,\qquad \operatorname{Var}(p_t)=\sigma^2t
\]

Variables: if \(p_0=0\), the note's simplified mean is \(\mu t\).

\subsubsection{Unit-root AR form}\label{unit-root-ar-form}

\[
y_t=\phi_1y_{t-1}+\text{error}_t
\]

Variables: \(\phi_1\) is the autoregressive coefficient.

\subsubsection{Unit root condition}\label{unit-root-condition}

\[
\phi_1=1
\]

Variables: a unit root implies random-walk-type non-stationarity.

\subsubsection{ADF regression}\label{adf-regression}

\[
\Delta y_t=c_t+\phi_cy_{t-1}+\beta_1\Delta y_{t-1}+\cdots+\beta_p\Delta y_{t-p}+\epsilon_t
\]

Variables: \(\Delta y_t=y_t-y_{t-1}\); \(c_t\) is deterministic
component; \(\phi_c\) is the unit-root test coefficient; \(\beta_j\)
controls lagged differences.

\subsubsection{ADF coefficient
relationship}\label{adf-coefficient-relationship}

\[
\phi_c=\phi_1-1
\]

Variables: \(\phi_c=0\) corresponds to \(\phi_1=1\).

\subsubsection{ADF hypotheses}\label{adf-hypotheses}

\[
\begin{aligned}
H_0&:\phi_c=0\quad \text{unit root / non-stationary}\\
H_1&:\phi_c<0\quad \text{stationary}
\end{aligned}
\]

Variables: the ADF alternative is one-sided and negative.

\subsubsection{ADF statistic}\label{adf-statistic}

\[
\operatorname{ADF}=\frac{\hat\phi_c}{\operatorname{se}(\hat\phi_c)}
\]

Variables: \(\hat\phi_c\) is the estimated ADF coefficient.

\subsubsection{ADF decision rule}\label{adf-decision-rule}

\[
\operatorname{ADF}<\text{ADF critical value}\Rightarrow \text{reject }H_0
\]

Variables: ``less than'' means more negative than the critical value.

\subsubsection{Drift-only ADF deterministic
component}\label{drift-only-adf-deterministic-component}

\[
c_t=c
\]

Variables: \(c\) is a constant drift/intercept.

\subsubsection{Trend ADF deterministic
component}\label{trend-adf-deterministic-component}

\[
c_t=c_0+c_1t
\]

Variables: \(c_0\) is an intercept; \(c_1\) is deterministic trend
slope.

\subsubsection{Martingale difference
sequence}\label{martingale-difference-sequence}

\[
\mathbb E(\epsilon_t\mid F_{t-1})=0
\]

Variables: \(\epsilon_t\) has zero conditional mean given past
information.

\subsubsection{AR(1) model}\label{ar1-model}

\[
y_t=c+\phi_1y_{t-1}+\epsilon_t
\]

Variables: \(c\) is intercept; \(\phi_1\) is the AR(1) coefficient.

\subsubsection{AR(1) conditional mean}\label{ar1-conditional-mean}

\[
\mathbb E(y_t\mid F_{t-1})=c+\phi_1y_{t-1}
\]

Variables: the conditional mean is the one-step forecast given past
information.

\subsubsection{AR(1) unconditional mean}\label{ar1-unconditional-mean}

\[
\mathbb E(y_t)=\frac{c}{1-\phi_1}
\]

Variables: this requires stationarity, typically \(|\phi_1|<1\).

\subsubsection{AR(1) conditional
variance}\label{ar1-conditional-variance}

\[
\operatorname{Var}(\epsilon_t\mid F_{t-1})=\sigma_\epsilon^2
\]

Variables: \(\sigma_\epsilon^2\) is innovation variance.

\[
\operatorname{Var}(y_t\mid F_{t-1})=\sigma_\epsilon^2
\]

Variables: uncertainty in \(y_t\) conditional on \(F_{t-1}\) comes from
the new shock.

\subsubsection{AR(1) unconditional
variance}\label{ar1-unconditional-variance}

\[
\operatorname{Var}(y_t)=\frac{\sigma_\epsilon^2}{1-\phi_1^2}
\]

Variables: this requires stationarity.

\subsubsection{AR(1) one-step forecast}\label{ar1-one-step-forecast}

\[
\mathbb E(y_{t+1}\mid F_t)=c+\phi_1y_t
\]

Variables: \(y_t\) is known at time \(t\).

\subsubsection{AR(1) two-step forecast}\label{ar1-two-step-forecast}

\[
\mathbb E(y_{t+2}\mid F_t)=c+\phi_1\mathbb E(y_{t+1}\mid F_t)=c(1+\phi_1)+\phi_1^2y_t
\]

Variables: the unknown future \(y_{t+1}\) is replaced by its forecast.

\subsubsection{AR(1) long-horizon
forecast}\label{ar1-long-horizon-forecast}

\[
\mathbb E(y_{t+h}\mid F_t)\to \frac{c}{1-\phi_1}
\]

Variables: forecasts mean-revert to the unconditional mean when the
AR(1) is stationary.

\subsection{ARCH and Volatility
Models}\label{arch-and-volatility-models}

\subsubsection{White noise}\label{white-noise}

\[
\begin{aligned}
\mathbb E(u_t)&=0\\
\operatorname{Var}(u_t)&=\sigma^2\\
\operatorname{Corr}(u_t,u_{t-k})&=0,\quad k\ge1
\end{aligned}
\]

Variables: \(u_t\) is a white-noise shock.

\subsubsection{MDS conditional mean}\label{mds-conditional-mean}

\[
\mathbb E(u_t\mid F_{t-1})=0
\]

Variables: an MDS is unpredictable from past information.

\subsubsection{MDS implies zero unconditional
mean}\label{mds-implies-zero-unconditional-mean}

\[
\mathbb E(u_t)=\mathbb E[\mathbb E(u_t\mid F_{t-1})]=0
\]

Variables: this uses the law of iterated expectations.

\subsubsection{MDS implies no
autocovariance}\label{mds-implies-no-autocovariance}

\[
\operatorname{Cov}(u_t,u_{t-k})=0,\quad k\ge1
\]

Variables: \(u_{t-k}\) is known at time \(t-1\).

\subsubsection{AR(p) model}\label{arp-model}

\[
y_t=c+\phi_1y_{t-1}+\cdots+\phi_py_{t-p}+\epsilon_t
\]

Variables: \(p\) is the autoregressive lag order.

\subsubsection{AR(p) unconditional mean}\label{arp-unconditional-mean}

\[
\mathbb E(y_t)=\frac{c}{1-\phi_1-\cdots-\phi_p}
\]

Variables: denominator uses the sum of AR coefficients.

\subsubsection{AR(p) one-step forecast}\label{arp-one-step-forecast}

\[
\mathbb E(y_{t+1}\mid F_t)=c+\phi_1y_t+\phi_2y_{t-1}+\cdots+\phi_py_{t-p+1}
\]

Variables: all lagged \(y\) values on the right are known at time \(t\).

\subsubsection{AR(p) selection}\label{arp-selection}

\[
\text{lowest }\operatorname{AIC}=\text{preferred model}
\]

Variables: \(\operatorname{AIC}\) is Akaike information criterion.

\subsubsection{Mean equation with
residual}\label{mean-equation-with-residual}

\[
r_t=\mathbb E(r_t\mid F_{t-1})+\epsilon_t
\]

Variables: \(\epsilon_t\) is the return shock after using information at
\(t-1\).

\subsubsection{Constant-variance
benchmark}\label{constant-variance-benchmark}

\[
\operatorname{Var}(\epsilon_t\mid F_{t-1})=\sigma_\epsilon^2
\]

Variables: this benchmark has no time-varying conditional volatility.

\subsubsection{ARCH LM mean model}\label{arch-lm-mean-model}

\[
y_t=x_t'\beta+\epsilon_t
\]

Variables: \(x_t\) is a regressor vector; \(\beta\) is a coefficient
vector.

\subsubsection{ARCH LM auxiliary
regression}\label{arch-lm-auxiliary-regression}

\[
\hat\epsilon_t^2=\rho_0+\rho_1\hat\epsilon_{t-1}^2+\cdots+\rho_q\hat\epsilon_{t-q}^2+v_t
\]

Variables: \(\hat\epsilon_t^2\) is squared residual; \(q\) is the number
of ARCH lags tested.

\subsubsection{ARCH LM hypotheses and
statistic}\label{arch-lm-hypotheses-and-statistic}

\[
\begin{aligned}
H_0&:\rho_1=\cdots=\rho_q=0\quad \text{no ARCH effects}\\
H_1&:\text{at least one }\rho_j\ne0\quad \text{ARCH effects exist}\\
TR^2&\sim\chi_q^2
\end{aligned}
\]

Variables: \(R^2\) comes from the auxiliary regression.

\subsubsection{ARCH LM decision rule}\label{arch-lm-decision-rule}

\[
TR^2>\chi_q^2\text{ critical value}\Rightarrow \text{reject }H_0
\]

Variables: rejecting means volatility is time-varying and predictable.

\subsubsection{ARCH model structure}\label{arch-model-structure}

\[
\epsilon_t=\sigma_tu_t,\qquad u_t\sim\mathrm{iid}(0,1)
\]

Variables: \(\sigma_t\) is conditional standard deviation; \(u_t\) is a
standardised shock.

\subsubsection{ARCH conditional moments}\label{arch-conditional-moments}

\[
\mathbb E(\epsilon_t\mid F_{t-1})=0,\qquad \operatorname{Var}(\epsilon_t\mid F_{t-1})=\sigma_t^2
\]

Variables: \(\sigma_t^2\) is conditional variance.

\subsubsection{ARCH(1)}\label{arch1}

\[
\sigma_t^2=\alpha_0+\alpha_1\epsilon_{t-1}^2
\]

Variables: \(\alpha_0\) is variance intercept; \(\alpha_1\) is the ARCH
coefficient.

\subsubsection{ARCH(1) conditions}\label{arch1-conditions}

\[
\alpha_0>0,\qquad \alpha_1\ge0,\qquad \alpha_1<1
\]

Variables: these keep variance positive and finite.

\subsubsection{ARCH(q)}\label{archq}

\[
\sigma_t^2=\alpha_0+\alpha_1\epsilon_{t-1}^2+\cdots+\alpha_q\epsilon_{t-q}^2
\]

Variables: \(q\) is the number of lagged squared shocks.

\subsubsection{ARCH(q) conditions}\label{archq-conditions}

\[
\alpha_0>0,\qquad \alpha_i\ge0,\qquad \sum_{i=1}^{q}\alpha_i<1
\]

Variables: \(\sum\alpha_i\) measures volatility persistence in ARCH(q).

\subsubsection{ARCH(1) unconditional
moments}\label{arch1-unconditional-moments}

\[
\begin{aligned}
\mathbb E(\epsilon_t)&=0\\
\operatorname{Cov}(\epsilon_t,\epsilon_{t-k})&=0,\quad k\ge1\\
\operatorname{Var}(\epsilon_t)&=\frac{\alpha_0}{1-\alpha_1}
\end{aligned}
\]

Variables: ARCH errors are uncorrelated, but squared errors can be
autocorrelated.

\subsubsection{Information criteria}\label{information-criteria}

\[
AIC=-2\ell+2k,\qquad BIC=-2\ell+k\log(T)
\]

Variables: \(\ell\) is log-likelihood; \(k\) is number of estimated
parameters; lower is better.

\subsection{GARCH, Forecasting Volatility, and Model
Comparison}\label{garch-forecasting-volatility-and-model-comparison}

\subsubsection{ARCH(q) variance equation}\label{archq-variance-equation}

\[
\sigma_t^2=\alpha_0+\alpha_1\epsilon_{t-1}^2+\cdots+\alpha_q\epsilon_{t-q}^2
\]

Variables: this is the pure ARCH variance equation.

\subsubsection{GARCH(p,q)}\label{garchpq}

\[
\begin{aligned}
\epsilon_t&=\sigma_tu_t,\qquad u_t\sim\mathrm{iid}(0,1)\\
\sigma_t^2&=\alpha_0+\alpha_1\epsilon_{t-1}^2+\cdots+\alpha_q\epsilon_{t-q}^2\\
&\quad+\beta_1\sigma_{t-1}^2+\cdots+\beta_p\sigma_{t-p}^2
\end{aligned}
\]

Variables: \(\beta_j\) are coefficients on lagged conditional variances;
\(p\) is GARCH order; \(q\) is ARCH order.

\subsubsection{GARCH(1,1)}\label{garch11}

\[
\sigma_t^2=\alpha_0+\alpha_1\epsilon_{t-1}^2+\beta_1\sigma_{t-1}^2
\]

Variables: \(\alpha_1\) is the shock reaction; \(\beta_1\) is volatility
persistence from lagged variance.

\subsubsection{GARCH(1,1) conditions}\label{garch11-conditions}

\[
\alpha_0>0,\qquad \alpha_1\ge0,\qquad \beta_1\ge0,\qquad \alpha_1+\beta_1<1
\]

Variables: \(\alpha_1+\beta_1<1\) gives finite unconditional variance
and mean reversion.

\subsubsection{GARCH persistence}\label{garch-persistence}

\[
\rho=\alpha_1+\beta_1
\]

Variables: \(\rho\) measures persistence of volatility shocks.

\subsubsection{GARCH long-run variance and
volatility}\label{garch-long-run-variance-and-volatility}

\[
\operatorname{Var}(\epsilon_t)=\bar\sigma^2=\frac{\alpha_0}{1-\alpha_1-\beta_1},\qquad
\bar\sigma=\sqrt{\frac{\alpha_0}{1-\alpha_1-\beta_1}}
\]

Variables: \(\bar\sigma^2\) is long-run variance; \(\bar\sigma\) is
long-run volatility.

\subsubsection{GARCH conditional and unconditional
moments}\label{garch-conditional-and-unconditional-moments}

\[
\begin{aligned}
\mathbb E(\epsilon_t\mid F_{t-1})&=0\\
\operatorname{Var}(\epsilon_t\mid F_{t-1})&=\sigma_t^2\\
\operatorname{Cov}(\epsilon_t,\epsilon_{t-k}\mid F_{t-1})&=0
\end{aligned}
\]

Variables: conditional variance changes over time even when conditional
mean is zero.

\[
\begin{aligned}
\mathbb E(\epsilon_t)&=0\\
\operatorname{Cov}(\epsilon_t,\epsilon_{t-k})&=0\\
\operatorname{Var}(\epsilon_t)&=\frac{\alpha_0}{1-\alpha_1-\beta_1}
\end{aligned}
\]

Variables: standard GARCH errors are uncorrelated but not independent in
squares.

\subsubsection{Normal GARCH innovation
skewness}\label{normal-garch-innovation-skewness}

\[
u_t\sim\mathrm{iid}\mathcal N(0,1),\qquad \epsilon_t\mid F_{t-1}\sim\mathcal N(0,\sigma_t^2)
\]

Variables: conditional normality makes conditional skewness zero.

\[
\mathbb E(\epsilon_t^3\mid F_{t-1})=0,\qquad
\mathbb E(\epsilon_t^3)=\mathbb E[\mathbb E(\epsilon_t^3\mid F_{t-1})]=0
\]

Variables: the unconditional third moment is zero under symmetric normal
innovations.

\[
\text{skewness}=\frac{\mathbb E(\epsilon_t^3)}{[\operatorname{Var}(\epsilon_t)]^{3/2}}=0
\]

Variables: skewness is zero for standard symmetric GARCH innovations.

\subsubsection{Mean equation plus GARCH volatility
equation}\label{mean-equation-plus-garch-volatility-equation}

\[
\begin{aligned}
r_t&=c+\phi_1r_{t-1}+\epsilon_t\\
\sigma_t^2&=\alpha_0+\alpha_1\epsilon_{t-1}^2+\beta_1\sigma_{t-1}^2
\end{aligned}
\]

Variables: \(c,\phi_1\) define the mean equation; the second line
defines volatility.

\subsubsection{Constant-mean return
equation}\label{constant-mean-return-equation}

\[
r_t=\mu+\epsilon_t
\]

Variables: \(\mu\) is constant mean return.

\subsubsection{Mean-adjusted AR form}\label{mean-adjusted-ar-form}

\[
r_t-\mu=\phi(r_{t-1}-\mu)+\epsilon_t
\]

Variables: \(\phi\) is the AR coefficient around mean \(\mu\).

\[
r_t=c+\phi r_{t-1}+\epsilon_t,\qquad c=\mu(1-\phi)
\]

Variables: \(c\) is the equivalent intercept.

\[
r_t=\mu+\phi(r_{t-1}-\mu)+\epsilon_t=\mu(1-\phi)+\phi r_{t-1}+\epsilon_t
\]

Variables: these are algebraically equivalent AR(1) mean forms.

\subsubsection{AR(1)-GARCH(1,1) one-step mean
forecast}\label{ar1-garch11-one-step-mean-forecast}

\[
\mathbb E(r_{T+1}\mid F_T)=\hat r_{T+1}=c+\phi r_T
\]

Variables: \(r_T\) is the latest known return.

\subsubsection{GARCH one-step variance
forecast}\label{garch-one-step-variance-forecast}

\[
\hat\sigma_{T+1}^2=\mathbb E(\sigma_{T+1}^2\mid F_T)=\alpha_0+\alpha_1\epsilon_T^2+\beta_1\sigma_T^2
\]

Variables: \(\epsilon_T^2\) is the latest squared shock; \(\sigma_T^2\)
is the latest conditional variance.

\subsubsection{Prediction interval using forecast
variance}\label{prediction-interval-using-forecast-variance}

\[
r_{T+1}\mid F_T\sim\mathcal N(\text{mean forecast},\text{variance forecast})
\]

Variables: mean forecast is \(\hat r_{T+1}\); variance forecast is
\(\hat\sigma_{T+1}^2\).

\[
\text{mean forecast}\pm1.96\sqrt{\text{variance forecast}}
\]

Variables: \(1.96\) is the approximate 95 percent normal cutoff.

\subsubsection{Likelihood-ratio comparison of ARCH and
GARCH}\label{likelihood-ratio-comparison-of-arch-and-garch}

\[
LR=2(LL_{GARCH}-LL_{ARCH})=2(\ell_1-\ell_0)
\]

Variables: \(LL_{GARCH}\) is the unrestricted log likelihood;
\(LL_{ARCH}\) is the restricted log likelihood.

\[
LR\sim\chi_p^2
\]

Variables: \(p\) is the number of restrictions or extra parameters
tested.

\subsubsection{Model comparison rules}\label{model-comparison-rules}

\[
\begin{aligned}
\text{higher log likelihood}&\Rightarrow \text{better}\\
\text{lower AIC/BIC}&\Rightarrow \text{better}\\
\text{squared standardised residuals with no autocorrelation}&\Rightarrow \text{better diagnostics}
\end{aligned}
\]

Variables: these compare fitted volatility models.

\subsection{Asymmetric Volatility
Models}\label{asymmetric-volatility-models}

\subsubsection{Standard GARCH symmetry}\label{standard-garch-symmetry}

\[
\sigma_t^2=\alpha_0+\alpha_1\epsilon_{t-1}^2+\beta_1\sigma_{t-1}^2
\]

Variables: squaring removes the sign of the shock.

\[
(+a)^2=(-a)^2=a^2
\]

Variables: \(a\) is any shock size.

\subsubsection{GJR-GARCH model}\label{gjr-garch-model}

\[
\begin{aligned}
r_t&=\mu+\epsilon_t\\
\epsilon_t&=\sigma_tu_t\\
\sigma_t^2&=\alpha_0+\alpha_1\epsilon_{t-1}^2+\lambda I_{t-1}\epsilon_{t-1}^2+\beta_1\sigma_{t-1}^2
\end{aligned}
\]

Variables: \(\lambda\) is the leverage/asymmetry coefficient;
\(I_{t-1}\) is a bad-news indicator.

\subsubsection{GJR indicator}\label{gjr-indicator}

\[
I_{t-1}=
\begin{cases}
1,&\epsilon_{t-1}\le0\\
0,&\text{otherwise}
\end{cases}
\]

Variables: the indicator switches on for non-positive shocks.

\subsubsection{GJR good-news and bad-news
effects}\label{gjr-good-news-and-bad-news-effects}

\[
\begin{aligned}
\epsilon_{t-1}>0&:\quad \text{effect}=\alpha_1\epsilon_{t-1}^2\\
\epsilon_{t-1}\le0&:\quad \text{effect}=(\alpha_1+\lambda)\epsilon_{t-1}^2
\end{aligned}
\]

Variables: good news uses only \(\alpha_1\); bad news uses
\(\alpha_1+\lambda\).

\[
\text{good news effect}=\alpha_1,\qquad \text{bad news effect}=\alpha_1+\lambda
\]

Variables: a positive \(\lambda\) means bad news has a larger variance
effect in GJR.

\subsubsection{GJR leverage test}\label{gjr-leverage-test}

\[
\begin{aligned}
H_0&:\lambda=0\\
H_1&:\lambda\ne0\\
t&=\frac{\hat\lambda}{\operatorname{se}(\hat\lambda)}
\end{aligned}
\]

Variables: \(\hat\lambda\) is the estimated leverage coefficient.

\[
|t|>1.96\Rightarrow \text{reject }H_0\text{ at approximately 5 percent}
\]

Variables: rejecting means statistically significant asymmetry.

\subsubsection{GJR parameter conditions}\label{gjr-parameter-conditions}

\[
\alpha_0>0,\qquad \alpha_1\ge0,\qquad \beta_1\ge0,\qquad \alpha_1+\lambda\ge0
\]

Variables: these keep conditional variance non-negative for good and bad
news.

\subsubsection{GJR unconditional
variance}\label{gjr-unconditional-variance}

\[
\operatorname{Var}(\epsilon_t)=\frac{\alpha_0}{1-\alpha_1-\beta_1-\lambda/2}
\]

Variables: \(\lambda/2\) appears when positive and negative shocks are
equally likely.

\subsubsection{EGARCH model}\label{egarch-model}

\[
\log(\sigma_t^2)=\omega+\alpha_1u_{t-1}+\gamma_1(|u_{t-1}|-\mathbb E|u_{t-1}|)+\beta_1\log(\sigma_{t-1}^2)
\]

Variables: \(\omega\) is intercept; \(\alpha_1\) captures signed shock
asymmetry; \(\gamma_1\) captures magnitude effects; \(\beta_1\) captures
log-volatility persistence.

\subsubsection{EGARCH standardised
shock}\label{egarch-standardised-shock}

\[
u_{t-1}=\frac{\epsilon_{t-1}}{\sigma_{t-1}}
\]

Variables: \(u_{t-1}\) standardises the shock by conditional volatility.

\subsubsection{Expected absolute standard normal
shock}\label{expected-absolute-standard-normal-shock}

\[
\mathbb E|u_t|=\sqrt{\frac{2}{\pi}}\approx0.7979
\]

Variables: this value is used when \(u_t\) is standard normal.

\subsubsection{EGARCH positivity}\label{egarch-positivity}

\[
\log(\sigma_t^2)\text{ is modelled }\Rightarrow \sigma_t^2\text{ is automatically positive}
\]

Variables: modelling log variance avoids non-negativity constraints on
coefficients.

\subsubsection{EGARCH leverage sign in the
notes}\label{egarch-leverage-sign-in-the-notes}

\[
\lambda<0\Rightarrow \text{negative shocks raise volatility more than positive shocks}
\]

Variables: \(\lambda\) is the asymmetry coefficient in the note's EGARCH
interpretation.

\subsubsection{News impact curve for
ARCH(1)}\label{news-impact-curve-for-arch1}

\[
NIC(\epsilon_{t-1})=\sigma_t^2=\alpha_0+\alpha_1\epsilon_{t-1}^2
\]

Variables: \(NIC\) maps yesterday's shock into today's conditional
variance.

\subsubsection{News impact curve for
GARCH(1,1)}\label{news-impact-curve-for-garch11}

\[
\bar\sigma^2=\frac{\alpha_0}{1-\alpha_1-\beta_1}
\]

Variables: \(\bar\sigma^2\) is the variance level used to hold lagged
variance fixed.

\[
NIC(\epsilon_{t-1})=\alpha_0+\alpha_1\epsilon_{t-1}^2+\beta_1\bar\sigma^2
\]

Variables: the curve varies with \(\epsilon_{t-1}\) while holding
\(\sigma_{t-1}^2=\bar\sigma^2\).

\[
NIC(\epsilon_{t-1})=\alpha_0+\beta_1\frac{\alpha_0}{1-\alpha_1-\beta_1}+\alpha_1\epsilon_{t-1}^2
\]

Variables: this substitutes the long-run variance into the NIC.

\subsubsection{News impact curve for
GJR-GARCH}\label{news-impact-curve-for-gjr-garch}

\[
A=\alpha_0+\beta_1\frac{\alpha_0}{1-\alpha_1-\lambda/2-\beta_1}
\]

Variables: \(A\) is the intercept-like NIC component after fixing lagged
variance at the GJR long-run variance.

\[
\begin{aligned}
\epsilon_{t-1}>0&:\quad \sigma_t^2=A+\alpha_1\epsilon_{t-1}^2\\
\epsilon_{t-1}\le0&:\quad \sigma_t^2=A+(\alpha_1+\lambda)\epsilon_{t-1}^2
\end{aligned}
\]

Variables: the two branches show asymmetric curvature for good and bad
news.

\subsection{IGARCH, RiskMetrics, MA, and
ARMA}\label{igarch-riskmetrics-ma-and-arma}

\subsubsection{IGARCH persistence}\label{igarch-persistence}

\[
\alpha_1+\beta_1=1
\]

Variables: IGARCH has unit volatility persistence.

\subsubsection{IGARCH(1,1) from GARCH(1,1)}\label{igarch11-from-garch11}

\[
\sigma_t^2=\alpha_0+\alpha_1\epsilon_{t-1}^2+\beta_1\sigma_{t-1}^2,\qquad \alpha_1=1-\beta_1
\]

Variables: the second condition enforces \(\alpha_1+\beta_1=1\).

\subsubsection{IGARCH with zero
intercept}\label{igarch-with-zero-intercept}

\[
\alpha_0=0
\]

Variables: RiskMetrics/EWMA uses zero variance intercept.

\[
\sigma_t^2=(1-\beta_1)\epsilon_{t-1}^2+\beta_1\sigma_{t-1}^2
\]

Variables: \(\beta_1\) is the decay weight on old variance.

\subsubsection{IGARCH using returns as
shocks}\label{igarch-using-returns-as-shocks}

\[
\sigma_t^2=(1-\beta_1)r_{t-1}^2+\beta_1\sigma_{t-1}^2
\]

Variables: this uses \(r_{t-1}\) as the shock when the conditional mean
is zero.

\subsubsection{IGARCH unconditional variance
issue}\label{igarch-unconditional-variance-issue}

\[
\operatorname{Var}(\epsilon_t)=\frac{\alpha_0}{1-\alpha_1-\beta_1},\qquad \alpha_1+\beta_1=1
\]

Variables: the denominator is zero, so stationary unconditional variance
is not finite.

\subsubsection{RiskMetrics conditional normal
model}\label{riskmetrics-conditional-normal-model}

\[
r_t\mid F_{t-1}\sim\mathcal N(0,\sigma_t^2)
\]

Variables: RiskMetrics often assumes zero conditional mean.

\subsubsection{RiskMetrics variance
update}\label{riskmetrics-variance-update}

\[
\sigma_t^2=\lambda\sigma_{t-1}^2+(1-\lambda)r_{t-1}^2
\]

Variables: \(\lambda\) is the decay factor; \(1-\lambda\) is the latest
squared-return weight.

\subsubsection{RiskMetrics one-step
forecast}\label{riskmetrics-one-step-forecast}

\[
\sigma_{T+1}^2=\lambda\sigma_T^2+(1-\lambda)r_T^2
\]

Variables: \(r_T\) is the latest observed return.

\subsubsection{Conditional-normal VaR}\label{conditional-normal-var}

\[
r_{T+1}\mid F_T\sim\mathcal N(\mu_{T+1\mid T},\sigma_{T+1}^2)
\]

Variables: \(\mu_{T+1\mid T}\) is the conditional mean forecast.

\[
q_\alpha=\mu_{T+1\mid T}+\sigma_{T+1}z_\alpha
\]

Variables: \(q_\alpha\) is the conditional return quantile.

\[
\operatorname{VaR}_\alpha=|W_0q_\alpha|
\]

Variables: \(W_0\) is dollar exposure.

\[
q_{0.05}=-1.645\sigma_{T+1},\qquad \operatorname{VaR}_{0.05}=|W_0(-1.645\sigma_{T+1})|
\]

Variables: this zero-mean form uses \(z_{0.05}=-1.645\).

\subsubsection{AR-GARCH VaR model}\label{ar-garch-var-model}

\[
\begin{aligned}
r_t&=\mu_t+\epsilon_t\\
\mu_t&=c+\phi_1r_{t-1}+\cdots+\phi_pr_{t-p}\\
\sigma_t^2&=\alpha_0+\alpha_1\epsilon_{t-1}^2+\beta_1\sigma_{t-1}^2
\end{aligned}
\]

Variables: \(\mu_t\) is conditional mean; \(p\) is AR order.

\subsubsection{AR-GARCH mean forecast}\label{ar-garch-mean-forecast}

\[
\mu_{T+1\mid T}=c+\phi_1r_T+\cdots+\phi_pr_{T-p+1}
\]

Variables: all returns on the right are known at time \(T\).

\subsubsection{AR-GARCH VaR quantile}\label{ar-garch-var-quantile}

\[
q_{0.05}=\mu_{T+1\mid T}-1.645\sigma_{T+1},\qquad \operatorname{VaR}_{0.05}=|W_0q_{0.05}|
\]

Variables: \(\sigma_{T+1}=\sqrt{\sigma_{T+1}^2}\).

\subsubsection{Persistence versus volatility
level}\label{persistence-versus-volatility-level}

\[
\text{persistence}=\alpha_1+\beta_1,\qquad \text{long-run variance}=\frac{\alpha_0}{1-\alpha_1-\beta_1}
\]

Variables: persistence controls shock decay; \(\alpha_0\) also affects
the volatility level.

\subsubsection{MA(1) model}\label{ma1-model}

\[
y_t=\phi_0+u_t+\theta_1u_{t-1}
\]

Variables: \(\phi_0\) is mean/intercept; \(\theta_1\) is MA coefficient;
\(u_t\) is shock.

\[
u_t\sim\mathrm{iid}(0,\sigma_u^2)
\]

Variables: \(\sigma_u^2\) is innovation variance.

\subsubsection{MA(1) mean, variance, and
autocorrelation}\label{ma1-mean-variance-and-autocorrelation}

\[
\mathbb E(y_t)=\phi_0,\qquad \operatorname{Var}(y_t)=(1+\theta_1^2)\sigma_u^2
\]

Variables: MA(1) variance depends on \(\theta_1^2\).

\[
\rho_1=\frac{\theta_1}{1+\theta_1^2},\qquad \rho_k=0\text{ for }k>1
\]

Variables: MA(1) autocorrelation cuts off after lag 1.

\subsubsection{MA(1) forecasts}\label{ma1-forecasts}

\[
y_{T+1}=\phi_0+u_{T+1}+\theta_1u_T
\]

Variables: \(u_T\) is the latest known/estimated shock.

\[
\hat y_{T+1}=\hat\phi_0+\hat\theta_1\hat u_T,\qquad \hat y_{T+h}=\hat\phi_0\text{ for }h\ge2
\]

Variables: future shocks have zero expected value.

\subsubsection{ARMA(p,q)}\label{armapq}

\[
y_t=\phi_0+\phi_1y_{t-1}+\cdots+\phi_py_{t-p}+u_t+\theta_1u_{t-1}+\cdots+\theta_qu_{t-q}
\]

Variables: \(p\) is AR order; \(q\) is MA order.

\[
\text{lowest }\operatorname{AIC}=\text{preferred}
\]

Variables: use AIC to choose among candidate ARMA models.

\subsection{Multi-Step Forecasting and Portfolio
VaR}\label{multi-step-forecasting-and-portfolio-var}

\subsubsection{Future squared shocks equal future variance in
expectation}\label{future-squared-shocks-equal-future-variance-in-expectation}

\[
\mathbb E(\epsilon_{t+h-1}^2\mid F_t)=\mathbb E(\sigma_{t+h-1}^2\mid F_t)
\]

Variables: future squared shocks are unknown, so use their conditional
expectation.

\[
\epsilon_{t+h-1}=\sigma_{t+h-1}u_{t+h-1},\qquad \mathbb E(u_{t+h-1}^2\mid F_{t+h-2})=1
\]

Variables: \(u\) is standardised to have conditional second moment 1.

\subsubsection{Recursive h-step GARCH variance
forecast}\label{recursive-h-step-garch-variance-forecast}

\[
\rho=\alpha_1+\beta_1
\]

Variables: \(\rho\) is GARCH persistence.

\[
\mathbb E(\sigma_{t+h}^2\mid F_t)=\alpha_0+\rho\mathbb E(\sigma_{t+h-1}^2\mid F_t)
\]

Variables: the \(h\)-step forecast uses the previous horizon's variance
forecast.

\subsubsection{Closed-form h-step GARCH
forecast}\label{closed-form-h-step-garch-forecast}

\[
\begin{aligned}
\mathbb E(\sigma_{t+h}^2\mid F_t)
&=\alpha_0\frac{1-\rho^{h-1}}{1-\rho}\\
&\quad+\rho^{h-1}\mathbb E(\sigma_{t+1}^2\mid F_t)
\end{aligned}
\]

Variables: \(h\) is forecast horizon; \(\rho^{h-1}\) controls mean
reversion speed.

\subsubsection{Closed-form forecast using long-run
variance}\label{closed-form-forecast-using-long-run-variance}

\[
\bar\sigma^2=\frac{\alpha_0}{1-\alpha_1-\beta_1}
\]

Variables: \(\bar\sigma^2\) is long-run variance.

\[
\mathbb E(\sigma_{t+h}^2\mid F_t)=\bar\sigma^2+\rho^{h-1}\left[\mathbb E(\sigma_{t+1}^2\mid F_t)-\bar\sigma^2\right]
\]

Variables: forecasts converge to \(\bar\sigma^2\) when \(\rho<1\).

\[
\mathbb E(\sigma_{t+h}^2\mid F_t)\to\bar\sigma^2
\]

Variables: this is the long-horizon mean-reversion result.

\subsubsection{Prediction interval for
returns}\label{prediction-interval-for-returns}

\[
r_{T+1}\mid F_T\sim\mathcal N(\mu_{T+1\mid T},\sigma_{T+1}^2)
\]

Variables: \(\mu_{T+1\mid T}\) is the one-step mean forecast;
\(\sigma_{T+1}^2\) is the one-step variance forecast.

\[
\mu_{T+1\mid T}\pm1.96\sqrt{\sigma_{T+1}^2}
\]

Variables: this is an approximate 95 percent prediction interval.

\subsubsection{One-asset VaR with forecast
variance}\label{one-asset-var-with-forecast-variance}

\[
q_{0.05}=\mu_{T+1\mid T}-1.645\sqrt{\sigma_{T+1}^2}
\]

Variables: \(q_{0.05}\) is the 5 percent bad-return quantile.

\[
\operatorname{VaR}_{0.05}=|W_0q_{0.05}|
\]

Variables: \(W_0\) is position value.

\[
\operatorname{VaR}_{0.05}=|W_0(-1.645\sqrt{\sigma_{T+1}^2})|
\]

Variables: this is the zero-mean special case.

\subsubsection{Asset-specific RiskMetrics
forecast}\label{asset-specific-riskmetrics-forecast}

\[
\sigma_{i,t+1}^2=(1-\lambda_i)r_{i,t}^2+\lambda_i\sigma_{i,t}^2
\]

Variables: \(i\) indexes the asset; \(\lambda_i\) is asset-specific
decay factor.

\subsubsection{Asset-specific RiskMetrics
VaR}\label{asset-specific-riskmetrics-var}

\[
\operatorname{VaR}_i=\left|W_i(-1.645)\sqrt{\sigma_{i,t+1}^2}\right|
\]

Variables: \(W_i\) is the dollar exposure to asset \(i\).

\subsubsection{Two-asset portfolio return for
VaR}\label{two-asset-portfolio-return-for-var}

\[
r_{p,t}=wr_{m,t}+(1-w)r_{c,t}
\]

Variables: \(w\) is the portfolio weight in asset \(m\); \(1-w\) is the
weight in asset \(c\).

\[
w=\frac{\text{amount in asset }m}{\text{total portfolio value}},\qquad
1-w=\frac{\text{amount in asset }c}{\text{total portfolio value}}
\]

Variables: weights are based on dollar values.

\subsubsection{Two-asset portfolio variance for
VaR}\label{two-asset-portfolio-variance-for-var}

\[
\sigma_p^2=w^2\sigma_m^2+(1-w)^2\sigma_c^2+2w(1-w)\rho\sigma_m\sigma_c
\]

Variables: \(\rho\) is correlation between the two asset returns.

\subsubsection{Portfolio VaR from portfolio
volatility}\label{portfolio-var-from-portfolio-volatility}

\[
\operatorname{VaR}_p=\left|\text{total value}\times(-1.645)\times\sigma_p\right|
\]

Variables: \(\sigma_p=\sqrt{\sigma_p^2}\) is portfolio standard
deviation.

\subsubsection{Portfolio VaR from individual
VaRs}\label{portfolio-var-from-individual-vars}

\[
\operatorname{VaR}_p=\sqrt{\operatorname{VaR}_m^2+\operatorname{VaR}_c^2+2\rho\operatorname{VaR}_m\operatorname{VaR}_c}
\]

Variables: \(\operatorname{VaR}_m,\operatorname{VaR}_c\) are stand-alone
asset VaRs.

\subsubsection{Generic two-asset portfolio variance and VaR
notation}\label{generic-two-asset-portfolio-variance-and-var-notation}

\[
\sigma_p^2=w_A^2\sigma_A^2+w_B^2\sigma_B^2+2w_Aw_B\rho_{AB}\sigma_A\sigma_B
\]

Variables: \(A,B\) are generic assets.

\[
\operatorname{VaR}_p=W_0(1.645)\sigma_p
\]

Variables: this writes VaR as a positive loss number when mean return is
zero.

\[
\operatorname{VaR}_p=\sqrt{\operatorname{VaR}_A^2+\operatorname{VaR}_B^2+2\rho_{AB}\operatorname{VaR}_A\operatorname{VaR}_B}
\]

Variables: this is the same normal portfolio VaR logic expressed using
individual VaRs.

\subsection{Essential Exam Answer
Templates}\label{essential-exam-answer-templates}

\subsubsection{ADF template}\label{adf-template}

\[
\begin{aligned}
H_0&:\text{unit root / non-stationary}\\
H_1&:\text{stationary}\\
\operatorname{ADF}&=\frac{\text{coefficient}}{\text{standard error}}\\
\text{reject only if ADF is more negative than the critical value}
\end{aligned}
\]

Variables: the coefficient is the ADF coefficient on lagged level.

\subsubsection{General t-test template}\label{general-t-test-template}

\[
\begin{aligned}
H_0&:\beta=\beta_0\\
H_1&:\beta\ne\beta_0\\
t&=\frac{\hat\beta-\beta_0}{\operatorname{se}(\hat\beta)}\\
\text{compare }|t|\text{ to the critical value}
\end{aligned}
\]

Variables: \(\beta_0\) is the hypothesised value.

\subsubsection{CAPM versus multi-factor
template}\label{capm-versus-multi-factor-template}

\[
\begin{aligned}
H_0&:\text{all extra factor }\beta\text{s}=0\\
H_1&:\text{at least one extra }\beta\ne0\\
J&=\frac{RSS_0-RSS_1}{RSS_1/(T-K-1)}\\
J&>\chi^2_{\text{restrictions}}\Rightarrow\text{reject }H_0
\end{aligned}
\]

Variables: \(RSS_0\) is restricted-model RSS; \(RSS_1\) is
unrestricted-model RSS.

\subsubsection{ARCH test template}\label{arch-test-template}

\[
\begin{aligned}
&\text{estimate mean model}\\
&\text{square residuals}\\
&\text{regress squared residuals on }q\text{ lags}\\
&TR^2\sim\chi_q^2\\
&\text{reject means ARCH effects / time-varying volatility}
\end{aligned}
\]

Variables: \(q\) is number of lagged squared residuals in the auxiliary
regression.

\subsubsection{One-step GARCH forecast
template}\label{one-step-garch-forecast-template}

\[
\begin{aligned}
\text{mean forecast}&=\text{fitted mean equation using known latest returns}\\
\text{variance forecast}&=\alpha_0+\alpha_1\epsilon_T^2+\beta_1\sigma_T^2
\end{aligned}
\]

Variables: the latest return, shock, and variance are known at time
\(T\).

\subsubsection{Prediction interval
template}\label{prediction-interval-template}

\[
\text{mean forecast}\pm1.96\sqrt{\text{variance forecast}}
\]

Variables: use 1.96 for an approximate 95 percent interval under
normality.

\subsubsection{5 percent VaR template}\label{percent-var-template}

\[
\begin{aligned}
q_{0.05}&=\text{mean forecast}-1.645\sqrt{\text{variance forecast}}\\
\operatorname{VaR}_{0.05}&=|\text{position value}\times q_{0.05}|
\end{aligned}
\]

Variables: 1.645 is the one-sided 5 percent normal cutoff.

\end{document}
